\documentclass[11pt]{article}
\usepackage[T1]{fontenc}
\usepackage[utf8]{inputenc}
\usepackage[french]{babel}
\usepackage[a4paper,margin=2.2cm]{geometry}
\usepackage{amsmath,amssymb}
\usepackage{xcolor}
\usepackage{enumitem}
\usepackage{titlesec}
\usepackage{fancyhdr}
\usepackage{tcolorbox}
\usepackage{hyperref}

\definecolor{lpp}{HTML}{0F52AC}
\definecolor{teal}{HTML}{00A37A}
\definecolor{grisfonce}{HTML}{333333}

\hypersetup{colorlinks=true, linkcolor=lpp, urlcolor=lpp}

\titleformat{\section}{\Large\bfseries\color{lpp}}{}{0pt}{}
\titleformat{\subsection}{\large\bfseries\color{grisfonce}}{}{0pt}{}
\setlist[itemize]{leftmargin=1.3em, itemsep=2pt, topsep=2pt}

\pagestyle{fancy}
\fancyhf{}
\lhead{\small\color{lpp}\textbf{adc\_cpp} — notes de présentation}
\rhead{\small Romain Despoullains}
\cfoot{\small\thepage}
\renewcommand{\headrulewidth}{0.4pt}

\newcommand{\slidetag}[2]{\subsection{\textcolor{lpp}{[#1]}\ \ #2}}
\newcommand{\dit}[1]{\par\smallskip\textbf{Ce que je dis :}\ #1}
\newcommand{\msg}[1]{\par\textit{\textcolor{teal}{Message clé :}}\ \textit{#1}}
\newcommand{\code}[1]{\texttt{\small #1}}

\newtcolorbox{qa}[1]{colback=lpp!4, colframe=lpp!55, boxrule=0.5pt,
  left=6pt, right=6pt, top=4pt, bottom=4pt, title={#1}, fonttitle=\bfseries\small}

\begin{document}

\begin{center}
  {\LARGE\bfseries\color{lpp} adc\_cpp : architecture d'un solveur AMR / MPI / GPU}\\[2pt]
  {\large Script de présentation et questions anticipées}\\[2pt]
  {\small Romain Despoullains — LPP}
\end{center}
\vspace{0.4em}
\hrule
\vspace{0.8em}

\noindent\textbf{Format.} 13 diapositives, viser $\sim$12 min de parole, $\sim$1 min par diapo
(la timeline « construire son solveur » un peu plus). Garder le ton : on présente une
\emph{conception}, pas un catalogue de fonctionnalités. Le fil : une équation $\rightarrow$
deux contrats $\rightarrow$ un coupleur $\rightarrow$ on construit un solveur en 5 étapes.

\section{Script diapo par diapo}

\slidetag{Titre}{Architecture d'un solveur AMR / MPI / GPU}
\dit{« Je présente \code{adc\_cpp}, un solveur C++23 pour des systèmes couplés
hyperbolique–elliptique, du type qu'on rencontre en physique des plasmas. L'angle
n'est pas la physique d'un cas précis mais la \emph{conception} : comment on écrit un
code qui tourne en série, en OpenMP, en MPI et sur GPU sans réécrire les opérateurs. »}
\msg{On parle d'architecture logicielle au service de la physique.}

\slidetag{1/11}{Le problème}
\dit{« Le problème générique : une partie hyperbolique, le champ conservé $U$ qui est
transporté, et une partie elliptique, un potentiel $\phi$ solution d'une équation de
Poisson $\Delta\phi = f(U)$. Les deux sont couplées \emph{à chaque pas} par la variable
auxiliaire $\mathrm{aux} = (\phi, \nabla\phi)$ : le potentiel entre dans le flux de la
partie hyperbolique. Le diocotron, l'Euler–Poisson auto-gravitant rentrent dans ce
moule. »}
\msg{Toute la suite découle de ce découpage hyperbolique / elliptique / couplage.}

\slidetag{2/11}{Architecture du solveur}
\dit{« Cinq couches, une responsabilité chacune. \textbf{PhysicalModel} : les lois
locales (flux, sources, fermetures), et rien d'autre — c'est device-callable.
\textbf{NumericalMethod} : reconstruction, flux numérique, opérateur elliptique, CL.
\textbf{DataLayout} : les conteneurs, MultiFab / BoxArray / Geometry / hiérarchie AMR.
\textbf{ExecutionBackend} : les \emph{seams}, les coutures où le parallélisme est isolé —
\code{for\_each\_cell}, la comm MPI, l'allocateur. Au-dessus, \textbf{TimeIntegrator}
compose les opérateurs (RK, IMEX, splitting) sans connaître leur implémentation.
Le point clé : on n'écrit de la physique qu'au sommet, et le backend ne voit que des
vues minimales (une Box, un Array4, un scalaire, un rang). Il ne touche jamais à
BoxArray. C'est ça qui rend le passage GPU / MPI mécanique. »}
\msg{Les couches sont étanches ; le parallélisme vit dans une seule couche, les seams.}

\slidetag{3/11}{PhysicalModel (le concept)}
\dit{« Un modèle, c'est un \emph{type} qui fournit quatre fonctions locales : \code{flux},
\code{max\_wave\_speed} (qui donne le pas de temps CFL), \code{source}, et
\code{elliptic\_rhs} (le terme source de Poisson, le couplage). C'est un \code{concept}
C++20 : le compilateur vérifie le contrat. Pas d'héritage, pas de fonction virtuelle.
Tout type qui remplit ces quatre signatures est un modèle valide, et le code est
entièrement inliné. »}
\msg{Le contrat est vérifié à la compilation, pas par une hiérarchie de classes.}

\slidetag{4/11}{Exemple : le diocotron}
\dit{« Voici un modèle concret qui remplit ce contrat : le diocotron. Une seule variable,
la densité électronique $n_e$. Le flux est le transport par la dérive $E\times B$ :
$v = (E\times B)/|B|^2$, et $F = n_e v$. Le terme elliptique, c'est la densité moins un
fond ionique neutralisant. Toutes les fonctions sont marquées \code{ADC\_HD} : la même
ligne de code compile sur CPU et sur GPU. C'est exactement les quatre fonctions du
concept, écrites pour un cas physique. »}
\msg{Le concept de la diapo précédente, instancié sur une vraie physique.}

\slidetag{5/11}{EllipticSolver (le concept)}
\dit{« Le jumeau elliptique. Là où PhysicalModel décrit la physique locale,
EllipticSolver décrit « tout ce qui sait résoudre $\Delta\phi = f$ ». Trois méthodes :
\code{rhs()} où on écrit le second membre $f$, \code{solve()} qui résout, \code{phi()}
où on lit la solution. Le \code{Coupler<Model, Elliptic>} est le seul à connaître les
deux contrats : il pousse $f$, résout, relit $\phi$, calcule $\nabla\phi$, le range dans
\code{aux}, et le redonne au flux. Trois backends remplissent ce concept et sont
interchangeables : multigrille géométrique, FFT spectral, FFT distribué MPI. »}
\msg{Deux contrats indépendants, reliés par un seul coupleur via $\mathrm{aux}=(\phi,\nabla\phi)$.}

\slidetag{6/11}{Construire son solveur avec adc}
\dit{« Maintenant, la partie pratique : comment un utilisateur construit son solveur. En
cinq étapes — définir le modèle, le faire tourner sur grille uniforme, choisir le schéma,
passer en AMR, choisir le backend. À droite, le diocotron tournant en AMR : on voit les
patchs raffinés (les rectangles) suivre la zone de fort gradient. »}
\msg{De zéro à un solveur AMR/GPU en cinq gestes, sans toucher au moteur.}

\slidetag{7/11}{Étape 1 : définir son modèle}
\dit{« On reprend tout depuis le début avec un modèle d'advection scalaire à vitesse
$(a,b)$. On remplit les quatre fonctions, et la dernière ligne, le \code{static\_assert},
demande au compilateur de vérifier que le type satisfait bien \code{PhysicalModel}. Si une
signature manque, ça ne compile pas. »}
\msg{Écrire un modèle = remplir un \code{struct} et laisser le compilateur valider.}

\slidetag{8/11}{Étape 2 : grille uniforme}
\dit{« On instancie un \code{Coupler<Diocotron, GeometricMG>} avec une géométrie et des
conditions limites. La boucle en temps appelle \code{advance}. Le coupleur ferme la boucle
$\phi \rightarrow \mathrm{aux} \rightarrow$ avance, automatiquement. Pour un modèle
purement hyperbolique sans Poisson, on appellerait directement \code{advance\_ssprk2}. »}
\msg{Le couplage Poisson–transport est encapsulé ; l'utilisateur n'écrit que la boucle.}

\slidetag{9/11}{Étape 3 : choisir le schéma}
\dit{« Le flux et la reconstruction sont des \emph{policies} template, passées en
paramètre, pas des \code{if} dans une boucle. \code{advance<VanLeer>} donne du MUSCL van
Leer. On a Rusanov, HLL, HLLC pour le flux ; NoSlope, Minmod, van Leer, MC, WENO5-Z pour
la reconstruction. Changer de schéma ne touche jamais le modèle, et comme c'est résolu à
la compilation, il n'y a aucun surcoût de branche. »}
\msg{Le schéma numérique est un paramètre, orthogonal à la physique.}

\slidetag{10/11}{Étape 4 : passer en AMR}
\dit{« Même interface, un coupleur de plus : \code{AmrCouplerMP}. On déclare les niveaux
(grossier + patchs fins), on donne un critère de raffinement (ici : densité au-dessus d'un
seuil), et la boucle appelle \code{regrid} périodiquement puis \code{step}. Le moteur gère
la hiérarchie, le clustering Berger–Rigoutsos, le sous-cyclage et le reflux conservatif.
L'utilisateur ne fournit que le critère. »}
\msg{L'AMR est un changement de coupleur, pas une réécriture du solveur.}
\par\smallskip\footnotesize\textcolor{teal}{(Note : le footer de cette diapo dit « regrid,
reux » — coquille pour « reflux ». Dire « reflux » à l'oral.)}\normalsize

\slidetag{11/11}{Étape 5 : choisir le backend}
\dit{« Et le code source ne change pas. Le backend se choisit une seule fois, à
\code{cmake} : un drapeau pour OpenMP, un pour MPI, un pour Kokkos/GPU. Le seam
\code{for\_each\_cell} se compile en boucle série, en \code{\#pragma omp}, ou en
\code{Kokkos::parallel\_for} selon le drapeau. Passer en GPU, c'est une option de
compilation, pas du code. Et l'OpenMP est garanti bit-à-bit identique à la série. »}
\msg{Série, OpenMP, MPI, GPU : un drapeau cmake, opérateurs inchangés.}

\slidetag{Fin}{Thank you for any feedback}
\dit{« Pour résumer : une équation générique, deux contrats (physique locale, solveur
elliptique), un coupleur qui les relie, et un backend choisi à la compilation. On obtient
un solveur qui passe du laptop au GPU sans réécriture. Merci, je prends vos questions. »}

\section{Questions probables et réponses}

\subsection{Conception / C++}

\begin{qa}{Pourquoi des concepts et pas de l'héritage virtuel ?}
Performance et portabilité GPU. Une fonction virtuelle implique une vtable et un appel
indirect : ça casse l'inlining dans les boucles chaudes, et ce n'est pas appelable
proprement depuis un kernel device. Les concepts donnent un polymorphisme à la
compilation : tout est inliné, zéro surcoût, et le même code part sur GPU. Le compilateur
valide le contrat (\code{static\_assert}) au lieu d'une hiérarchie de classes.
\end{qa}

\begin{qa}{Le choix du backend à la compilation, ça ne manque pas de souplesse ?}
C'est un choix assumé : on cible le HPC, où la machine est connue au build (un nœud GPU,
un cluster MPI). Le gain est que les opérateurs restent monomorphes et pleinement
optimisés. La flexibilité « runtime » qu'on veut vraiment (résolution, schéma au niveau
utilisateur, modèle) reste accessible, notamment depuis Python qui pilote le binaire
compilé sans recompiler.
\end{qa}

\begin{qa}{Que veut dire \code{ADC\_HD} exactement ?}
C'est une macro : \code{\_\_host\_\_ \_\_device\_\_} quand on compile avec Kokkos/CUDA,
vide sinon. Elle marque les fonctions appelables à la fois sur l'hôte et sur le GPU. Comme
le modèle n'utilise que de l'arithmétique sur des petites vues (pas d'allocation, pas de
\code{std::} non-device), la même fonction compile pour les deux cibles.
\end{qa}

\begin{qa}{\code{StateVec<1>}, \code{n\_vars} : c'est quoi ?}
\code{State} est le vecteur d'état par cellule ; \code{StateVec<1>} = une variable
(le diocotron transporte juste $n_e$). Pour Euler 2D ce serait \code{StateVec<4>}
(densité, deux quantités de mouvement, énergie). \code{n\_vars} expose cette taille au
reste du code.
\end{qa}

\subsection{Physique}

\begin{qa}{C'est quoi l'instabilité diocotron ?}
L'instabilité d'une colonne de plasma non neutre (un faisceau d'électrons) en champs
croisés $E\times B$. C'est une instabilité de cisaillement, l'analogue électrostatique de
Kelvin–Helmholtz : le gradient de la vitesse de dérive déstabilise des perturbations
azimutales, qui croissent exponentiellement. Le taux de croissance dépend du nombre de
mode $l$.
\end{qa}

\begin{qa}{Comment $\nabla\phi$ entre dans le flux ?}
Le champ électrique est $E = -\nabla\phi$, et la vitesse de dérive est
$v = (E\times B)/|B|^2$. Dans le code, $v_x = -\partial_y\phi / B_0$ et
$v_y = +\partial_x\phi / B_0$. Le coupleur calcule $\nabla\phi$ par différences finies
centrées sur la grille après chaque résolution de Poisson, et le range dans \code{aux}.
\end{qa}

\begin{qa}{Le fond neutralisant $n_{i0}$ / la moyenne $\langle\rho\rangle$ : qui le calcule ?}
Le terme elliptique est $\alpha(n_e - n_{i0})$ : la densité moins un fond ionique. Pour un
Poisson périodique, le second membre doit être de moyenne nulle (nullspace). La quantité
globale ($\langle\rho\rangle$, le fond neutralisant) est calculée par le \emph{coupleur},
pas par le modèle local : le modèle ne voit qu'une cellule, il ne peut pas connaître une
moyenne globale. C'est un choix de conception : les quantités non locales remontent au
coupleur.
\end{qa}

\subsection{Numérique}

\begin{qa}{À quoi sert \code{max\_wave\_speed} ?}
Au pas de temps. La condition CFL borne $\Delta t$ par $\mathrm{CFL}\cdot\Delta x / v_{\max}$ ;
\code{max\_wave\_speed} fournit la vitesse caractéristique maximale locale, le coupleur en
prend le max global pour fixer $\Delta t$.
\end{qa}

\begin{qa}{van Leer, WENO5 : quand utiliser quoi ?}
NoSlope (ordre 1) pour du robuste/diagnostic ; les limiteurs MUSCL (Minmod, van Leer, MC)
pour de l'ordre 2 TVD sans osciller aux chocs ; WENO5-Z quand on veut de l'ordre élevé sur
des structures lisses (taux de croissance, turbulence). Le choix est un paramètre template,
donc on compare sans rien réécrire.
\end{qa}

\begin{qa}{Comment Poisson est-il résolu ?}
Trois backends. Multigrille géométrique (V-cycles, général, gère un bord conducteur
embarqué) ; FFT spectral direct (périodique, exact, très rapide) ; FFT distribué pour le
MPI. Tous remplissent le concept EllipticSolver, donc le coupleur est indifférent au choix.
\end{qa}

\subsection{Parallélisme / AMR / HPC}

\begin{qa}{AMR : sous-cyclage et reflux, c'est conservatif ?}
Oui. Le niveau fin fait plusieurs petits pas pour un pas grossier (sous-cyclage). Aux
bords grossier/fin, les flux ne coïncident pas : un \code{FluxRegister} accumule la
différence et corrige les cellules grossières (reflux). La somme utilise une accumulation
déterministe (compensée), donc la masse est conservée à l'arrondi (dérive $\sim 10^{-13}$
à $10^{-15}$) et le résultat est bit-à-bit reproductible.
\end{qa}

\begin{qa}{Berger–Rigoutsos, c'est quoi ?}
L'algorithme standard pour regrouper les cellules taguées en un petit nombre de patchs
rectangulaires efficaces. Il calcule des signatures (histogrammes des tags) et coupe là où
le laplacien de la signature change de signe (un trou ou un coin). On évite ainsi un patch
unique trop gros ou une myriade de petits.
\end{qa}

\begin{qa}{Comment se fait la décomposition MPI ?}
Un \code{BoxArray} pave le domaine en boîtes ; un \code{DistributionMapping} assigne
chaque boîte à un rang ; un \code{MultiFab} est la collection distribuée des données.
L'échange de halos (\code{fill\_boundary}) passe par le seam \code{comm}. Le modèle et les
schémas ne savent pas qu'on est en MPI : ils ne voient qu'une Box locale.
\end{qa}

\begin{qa}{OpenMP / GPU bit-identique à la série, comment ?}
Les réductions (max d'onde, masse) utilisent des réducteurs déterministes
(somme/max à ordre fixe), pas une accumulation atomique dépendante de l'ordre des threads.
On a vérifié : série et MPI donnent des sommes identiques bit-à-bit, et un run GPU
(GH200 sur ROMEO) donne le même checksum que le CPU, avec \code{compute-sanitizer} à zéro
erreur.
\end{qa}

\begin{qa}{À quoi sert l'IMEX / l'Asymptotic Preserving dans TimeIntegrator ?}
À traiter les termes raides implicitement. Quand le couplage devient raide (limite
quasi-neutre, sources rapides), un schéma explicite imposerait un $\Delta t$ minuscule.
IMEX traite la partie raide en implicite, et un schéma AP reste consistant et stable dans
la limite singulière. TimeIntegrator les compose sans connaître l'implémentation des
opérateurs.
\end{qa}

\subsection{Validation / limites}

\begin{qa}{Comment savez-vous que c'est correct ?}
Plusieurs niveaux. Taux de croissance diocotron comparés à l'analytique (modes 3/4/5,
WENO5-Z + SSPRK3) : on retrouve les taux avec un léger sur-tir qui est \emph{plat} en
résolution. L'AMR suit l'uniforme à résolution effective égale, avec $\sim 40\%$ des
cellules. Conservation de la masse à l'arrondi. Et la reproductibilité bit-à-bit
série/MPI/GPU sert de test de non-régression fort.
\end{qa}

\begin{qa}{Le sur-tir sur le taux de croissance, vous l'expliquez ?}
C'est une question ouverte : le sur-tir dépend du mode $l$ et reste structurel après avoir
écarté plusieurs pistes (grille, observable, normalisation). La piste actuelle est le
traitement des bords de l'anneau sur une grille cartésienne. C'est cohérent en résolution,
donc ce n'est pas un défaut de convergence mais un biais de modélisation du bord.
\end{qa}

\begin{qa}{En quoi c'est différent d'AMReX ou d'un framework existant ?}
Ce n'est pas un concurrent généraliste. C'est une conception ciblée : les contrats
(PhysicalModel, EllipticSolver) et les seams sont volontairement minces, pour que le
couplage hyperbolique–elliptique et le passage multi-backend soient lisibles et vérifiés.
L'idée pédagogique : montrer qu'avec deux concepts et une couche de seams, on couvre
série / OpenMP / MPI / GPU et l'AMR.
\end{qa}

\begin{qa}{Qu'est-ce qui n'est pas encore fait ?}
Honnêtement : le bord d'anneau cartésien (cf. sur-tir) ; l'extension à plus de modèles
multi-variables couplés ; et un durcissement du multigrille près du bord conducteur
embarqué à haute résolution, où le V-cycle peut diverger (documenté). Le socle
architecture / parallélisme / AMR est en place et testé.
\end{qa}

\end{document}
